\chapter{Quick start}\label{ch:quick}

\section{Installation}

ContourRoots is plain MATLAB code; the core needs no additional toolbox.
Download \code{ContourRoots.mltbx} from the release page and open it to
install it as an Add-On, or download \code{ContourRoots.zip}, unzip it and
run \code{setup\_contourroots}, which adds the toolbox to the path for the
current session. From the Command Window:
\begin{lstlisting}
url = "https://github.com/flavioluiz/ContourRoots/releases/latest/download/ContourRoots.zip";
unzip(websave("ContourRoots.zip", url), "ContourRoots");
run("ContourRoots/setup_contourroots.m")
\end{lstlisting}

\section{Roots of a characteristic function}

The characteristic function $F(s)=s^2+s+1+(2s+3)e^{-s}$ has infinitely
many roots. Those in $-8<\Real s<2$, $-20<\Imag s<20$ are obtained with
\begin{lstlisting}
F = @(s) s.^2 + s + 1 + (2*s + 3).*exp(-s);
[r, info] = croots(F, [-8 2 -20 20], 'AssumeAnalytic', true);
\end{lstlisting}
The vector \code{r} contains seven roots, ordered by decreasing real part;
the first two, $0.36201\pm1.8195\ii$, lie in the right half-plane.
\code{info.status} is \code{'numerically\_complete'}. The option
\code{'AssumeAnalytic',true} is the user's statement that $F$ is analytic
on a neighborhood of the closed rectangle; it is justified here because
$F$ is built from polynomials and the exponential. Without it the search
is exploratory (Section~\ref{sec:exploratory}).

\section{Poles and zeros}

A transfer function is given as a pair of analytic functions:
\begin{lstlisting}
G = ndpair(@(s) sinh(s/2), @(s) sinh(s));
p = cpoles(G, [-1 1 -10 10], 'AssumeAnalytic', true);   % +-i*pi, +-3i*pi
z = czeros(G, [-1 1 -10 10], 'AssumeAnalytic', true);   % empty
cpzmap(G, [-1 1 -10 10], 'AssumeAnalytic', true)
\end{lstlisting}
The denominator vanishes at $\ii\pi k$ for every integer $k$; for even $k$
the numerator vanishes as well and the candidate cancels.
Figure~\ref{fig:overview} shows both results.

\begin{figure}[H]
\centering
\includegraphics[width=\textwidth]{readme_overview.png}
\caption{Left: roots of $s^2+s+1+(2s+3)e^{-s}$ in the rectangle
$[-8,2]\times[-20,20]$ (dashed); the shaded area is the right half-plane.
Right: \code{cpzmap} of $\sinh(s/2)/\sinh(s)$; grey dots are cancelled
candidates.}
\label{fig:overview}
\end{figure}

\section{Time-delay stability}

For $D(s)+N(s)e^{-sT}=0$ with coefficient vectors \code{N} and \code{D}:
\begin{lstlisting}
crit = critical_delays(2, [1 1], 10);         % s + 1 + 2 e^{-sT} = 0
Z = unstable_root_count(2, [1 1], 1.5);       % 2 unstable roots
\end{lstlisting}
The first critical delay is $2\pi/(3\sqrt3)\approx1.2092$
(Chapter~\ref{ch:delay}).
